\documentclass[11pt]{article}
\usepackage[T1]{fontenc}
\usepackage{lmodern,amsmath,amssymb,amsthm,mathtools,booktabs,array,enumitem,microtype}
\usepackage[margin=0.9in]{geometry}
\usepackage{xcolor}
\usepackage[colorlinks=true,linkcolor=blue!55!black,urlcolor=blue!55!black,citecolor=blue!55!black]{hyperref}
\newtheorem{theorem}{Theorem}[section]
\newtheorem{lemma}[theorem]{Lemma}
\newtheorem{corollary}[theorem]{Corollary}
\newtheorem{proposition}[theorem]{Proposition}
\theoremstyle{definition}\newtheorem{definition}[theorem]{Definition}
\newcommand{\F}{\mathbb F}\newcommand{\ord}{\operatorname{ord}}\newcommand{\Nrm}{\operatorname{N}}\newcommand{\vpp}{v_p}
\newcommand{\src}[2]{\href{https://github.com/VladimirReshetnikov/ProveIt/blob/ee10f937fedac0955d28b7b860f4c44fab508204/Analysis/GeneralizedRiemannHypothesis/Lean/grh/src/#1}{\texttt{#2}}}
\newcommand{\prp}[2]{\href{https://github.com/VladimirReshetnikov/ProveIt/tree/ee10f937fedac0955d28b7b860f4c44fab508204/Analysis/GeneralizedRiemannHypothesis/docs/proposals/#1}{#2}}
\setlist{nosep,leftmargin=*}
\setlength{\emergencystretch}{3em}
\hypersetup{pdftitle={Three power tests instead of primitive labels},pdfsubject={Finite-field simplification of the theta-lift proof of Dirichlet GRH}}
\title{\vspace{-1.2em}\textbf{Three power tests instead of primitive labels}\\[4pt]\large A finite-field simplification of the theta-lift proof of Dirichlet GRH}
\author{Research proposal and integration blueprint}
\date{27 November 2026}
\begin{document}
\maketitle
\begin{abstract}
The readable Alp\"oge--Claude proof imposes a primitive-polynomial witness above each of two level primes. Primitivity serves two different downstream purposes: an inertia eigenvalue must retain its full Frobenius orbit after reduction at any coefficient prime, and a scalar inertia element must retain sufficiently large order. We replace the generator condition by three explicit power tests. For $q\equiv1\pmod6$ and $r=3^a>1$, the tests pick out exactly $(q^r-1)(1-1/C)/(3r)$ irreducible polynomials, where $C=(q^{2r/3}+q^{r/3}+1)/3$. They are invariant under the required tame scalings, yield the required integral matrix span, and provide the necessary scalar-order certificate with $6^s$ in place of the full tame order $q_2$. The witness count is at least $q^r/(8r)$, with no totient estimate on $q^r-1$. The resulting discarded Plancherel mass is at most $f(1+q_2+2r)/q+\exp(-f/(8r))$ per level prime. We keep all admissible scales and the arithmetic and analytic endpoint unchanged. The replacements below are mathematical proofs and a source-level integration plan, not a new Lean-checked RH proof.
\end{abstract}
\tableofcontents

\section{Location in the proof and scope}
The repository's readable development gives a one-sided estimate for the smoothed prime remainder by comparing the same twisted volume with local divisibility. The short-form and spanning theorems produce almost a basis whose theta lifts have a controlled covolume and index. The long-form theorem, from the determinant height and regular twists, gives the opposing bound. Cancelling the main term gives the source's \src{New/Reduction/MainInequality.lean\#L87}{star\_inequality},
\[
       m\,\mathrm{totalRemainder}\le f a\,(\log b)^5.
\]
At the analytic endpoint this yields the eventual upper bound $R_{q_0}(x)\le x^{-1/2+\epsilon}$, for the specified $q_0,d$. The one-sided criterion and the functional equations give the five statements in \src{Main.lean}{Main.lean}. See \src{New/Reduction/OneSidedBound.lean\#L200}{eventually\_primeRemainder\_le\_rpow} and the source's theorem \src{New/Reduction/OneSidedBound.lean\#L239}{grh}. The repository reports independent kernel checks of the original readable development; this note does not repeat them.

Here we change the local labels entering all three stages, not the volume inequalities, moment estimates, or analytic argument. At each of the two level primes $q=p_i$, the chosen local factors are copies of $\mathrm{GL}_r(\F_q)$. Their central projectors are indexed by monic squarefree degree-$r$ polynomials with nonzero constant term. Let $H=q_2$ denote the order of the fixed tame character (not the product of the two level primes, also denoted $M$ in the source). The existing good tuples satisfy:
\begin{enumerate}
\item every coordinate is allowed: no root in $\F_q$ of order dividing $H$;
\item at least one coordinate is primitive: irreducible with a root generating $\F_{q^r}^{\times}$.
\end{enumerate}
The first condition is essential for hyperplane annihilation and for the non-root condition on scaled polynomials. It will not be weakened. The second provides a cuspidal place for trace comparison and an inertial witness for the Selmer estimate. The same two-prime layout is retained, so for each coefficient prime $\ell$ at least one witness is at a different residue characteristic.

The key interface is \src{New/Congruences/GaloisRepresentations/RowsAndColumnsFromLabels.lean\#L326}{IsCommonRootPower}: after scaling, the roots have a common power $y$, are fixed by $q^r$-Frobenius, and have pairwise distinct $q$-conjugates even after taking any $\ell^u$-th powers. In addition, $y^x=1$ currently implies $H\mid jx$ for the lift index $j$. The scalar-separation argument uses this only to exclude $y^{\ell^u t}=1$ in a bounded range of $t$. In the actual data $H$ is far larger than necessary for this purpose. We exploit that observation to remove the global primitive-generator requirement.

\paragraph{Standing arithmetic assumptions.}
In the source, $r=3^a>1$, $r<q$, $q_0$ is divisible by $18480$, $s=\texttt{D.towerExp}$ satisfies $2^s=m^{22}$, and
\[
 H=(q_0d)^{s+6},\qquad H^2\mid q-1,\qquad
 0<m(4r+2)<2^s.
\]
Here $d$ is the auxiliary quadratic discriminant parameter. In particular $s\ge1$, $6^s\mid H$, $6\mid H$, and $q\equiv1\pmod6$. The last size inequality is precisely the bound used in the original \texttt{not\_q\textsubscript{2}\_dvd\_mul\_pow\_mul} in \src{New/Congruences/GaloisRepresentations/RowsAndColumnsFromLabels.lean\#L112}{RowsAndColumnsFromLabels.lean}. All these assumptions are already present; no new admissible-data existence theorem is required. In what follows the exponent $a$ in $r=3^a$ is unrelated to the derivative order $a$ in the main inequality.

\section{A replacement for full primitivity}
Set
\[
 A=q^{r/3},\quad B=A^2+A+1,\quad C=B/3,\quad N=q^r-1=(A-1)3C,
 \qquad E=N/(q-1).
\]
Since $A\equiv1\pmod3$, writing $A=1+3k$ gives $B=3+9k+9k^2$. Hence $v_3(B)=1$ and $3\nmid C$. Since $q$ is odd, $C$ is odd. Moreover $B\equiv3\pmod{q-1}$; thus $\gcd(C,q-1)=1$. We have $C>1$ and consequently $C\ge5$. In fact $C\ge19$ already for $q\ge7$ and $r\ge3$.

\begin{definition}[Robust label]
A robust element $z\in\F_{q^r}^\times$ satisfies
\begin{equation}\label{eq:tests}
 \boxed{\quad z^{N/2}\ne1,\qquad z^{N/3}\ne1,\qquad z^{3(A-1)}\ne1.\quad}
\end{equation}
A robust polynomial is a monic irreducible degree-$r$ polynomial over $\F_q$ with a root satisfying \eqref{eq:tests}. Its roots are nonzero and distinct, so it is a label. Equivalently, compute the three indicated powers of $X$ modulo the polynomial and check that none is the class of $1$.
\end{definition}

The tests are not an implicit primality or generator test. They use three exponents obtained by ordinary arithmetic, regardless of the number of prime factors of $N$. For $g$ a generator of the cyclic group of order $N$ and $z=g^h$, they are exactly
\begin{equation}\label{eq:hexp}
                 \gcd(h,6)=1,\qquad C\nmid h.
\end{equation}
This equivalence follows directly from the order of $g$. It is independent of the choice of generator, because changing generators multiplies $h$ by a unit modulo $N$. The formulation in \eqref{eq:tests} does not require choosing a generator when checking a polynomial.

\begin{lemma}[Exact orbit-erasure criterion]\label{lem:erasure}
For $z=g^h\in\F_{q^r}^\times$, the following are equivalent:
\begin{enumerate}[label=(\roman*)]
\item for every prime $\ell$ and every $u\ge0$, the $q$-Frobenius orbit of $z^{\ell^u}$ has size $r$;
\item $B/\gcd(B,h)$ has at least two distinct prime factors.
\end{enumerate}
In particular every robust element has this property. It is enough to require this at primes $\ell\ne\operatorname{char}\F_q$, but the stronger statement holds.
\end{lemma}
\begin{proof}
The Frobenius orbit size divides $r=3^a$. Its size is strictly smaller than $r$ exactly when $z^{\ell^u}$ is fixed by $q^{r/3}$-Frobenius, that is, exactly when
\[
 z^{(A-1)\ell^u}=1\quad\Longleftrightarrow\quad B\mid h\ell^u.
\]
The latter is equivalent to $B/\gcd(B,h)\mid \ell^u$. There exists a prime and a nonnegative power with this property if and only if $B/\gcd(B,h)$ is $1$ or a prime power. This proves the equivalence. Under \eqref{eq:hexp}, $3\nmid h$ and $C\nmid h$, so $B/\gcd(B,h)$ contains $3$ and a prime different from $3$.
\end{proof}

In particular a robust element has degree $r$ over $\F_q$ and yields a unique monic irreducible minimal polynomial of that degree. Every Frobenius conjugate is robust: $q$ is coprime to $N$, so $h\mapsto qh\pmod N$ preserves \eqref{eq:hexp}. The exact criterion is slightly weaker than our three tests: if $3\mid h$, the quotient in Lemma~\ref{lem:erasure} can still have two other prime factors. The stronger test $3\nmid h$ is intentionally retained because it also gives a useful scalar order.

\begin{lemma}[Tame scaling]\label{lem:scale}
Assume $6H\mid q-1$, and let $\alpha\in\F_q^\times$ satisfy $\alpha^H=1$. If $z$ is robust then $\alpha z$ is robust. In particular the source's stronger assumption $H^2\mid q-1$, $6\mid H$, suffices.
\end{lemma}
\begin{proof}
The exponents $N/2$ and $N/3$ are multiples of $H$, because $6H\mid q-1\mid N$. The exponent $3(A-1)$ is a multiple of $q-1$. Consequently $\alpha$ raised to any of the three exponents is $1$. Each test on $\alpha z$ has exactly the same value as on $z$. Multiplication by a nonzero element of the base field preserves degree and commutes with Frobenius, so the induced root scaling preserves the robust label condition.
\end{proof}

\begin{lemma}[Two-prime norm certificate]\label{lem:norm}
For a robust element $z$, put $y=z^E=\Nrm_{\F_{q^r}/\F_q}(z)$. Then
\[
 v_2(\ord(y))=v_2(q-1),\qquad v_3(\ord(y))=v_3(q-1).
\]
In particular $6^s\mid\ord(y)$ under the standing assumptions, and $z_i^E=y$ for every root $z_i$ of the robust polynomial. The same conclusion holds after any permitted tame scaling.
\end{lemma}
\begin{proof}
The norm lies in $\F_q^\times$, whose order is $q-1$. For $p=2,3$, $y^{(q-1)/p}=z^{N/p}\ne1$ by \eqref{eq:tests}. Since $\ord(y)$ divides $q-1$, failure to divide $(q-1)/p$ is equivalent to having its full $p$-part. The norm formula $E=1+q+\cdots+q^{r-1}$ shows that $z_i^E=y$ for all conjugates. Apply Lemma~\ref{lem:scale} to handle scaled roots. Finally $6^s\mid H\mid q-1$.
\end{proof}

\paragraph{A strictly larger class.}
For $q=37$, $r=3$, $H=6$, we have $N=50652$ and $C=469=7\cdot67$. Taking $h=7$ gives a robust element of order $7236$, not a generator. There are $5616$ robust polynomials and $864$ of them are not primitive. The class always contains the primitive polynomials. Nor must the norm itself be primitive: for $q=181$, $r=3$ and $h=5$ the element satisfies the tests and the norm has order $36$ rather than $180$. These are small finite-field models of the local hypotheses, not complete instances of the admissible theta-lift data.

\section{Why the Galois bridge still works}
\subsection{The precise weakened contract}
The scalar clause of \texttt{IsCommonRootPower} is the only clause which changes:
\begin{equation}\label{eq:contract}
             y^x=1\quad\Longrightarrow\quad 6^s\mid jx.
\end{equation}
All the other fields (common power, fixed by $q^r$, and injectivity after every $\ell^u$) are kept verbatim. At a robust full-rank label the scaled roots have common norm $y$ by Lemma~\ref{lem:norm}, so $6^s\mid x$, stronger than \eqref{eq:contract}. Lemma~\ref{lem:erasure} gives injectivity on $0,\ldots,r-1$ for each root. The proof uses the same tame scaling as before: at a chosen place the twist is unramified, and the scaling factor comes from the inverse of the tame character of order dividing $H$.

For the rank-one main critical pairs the old label is $X-1$ before scaling; do not impose the new robust tests here. Its scaled root $y$ has order $H/\gcd(H,j)$, since it is the name of $\omega^{-j}$ for a tame character of order $H$. Hence $y^x=1$ implies $H\mid jx$, and therefore \eqref{eq:contract}. The fixed-point and injectivity clauses in rank one are unchanged. This preserves both branches of the source's \texttt{exists\_isCommonRootPower}.

\begin{lemma}[Small-order obstruction]\label{lem:scalar}
Suppose $1\le j\le m$, $1\le t\le4r+2$ and $m(4r+2)<2^s$. For every prime $\ell$ and every $u\ge0$,
\[
                    6^s\nmid j\ell^u t.
\]
Consequently a root of unity $y$ satisfying \eqref{eq:contract} obeys $y^{\ell^u t}\ne1$ in this range.
\end{lemma}
\begin{proof}
Take $p=3$ if $\ell=2$, and $p=2$ otherwise. Then $p\ne\ell$. If $6^s\mid j\ell^u t$, then $p^s\mid jt$ by coprimality. But $0<jt\le m(4r+2)<2^s\le p^s$, a contradiction. The final assertion follows from \eqref{eq:contract}.
\end{proof}

For an injective multiplicative map $\tau$ from the algebraic closure of the finite residue field to the $\ell$-adic coefficient field, a root of unity has reduction $1$ if and only if its order is an $\ell$-power. Indeed its prime-to-$\ell$ component reduces injectively, and its $\ell$-primary component reduces to $1$. Thus the lemma gives exactly the unit bound needed by \src{New/Congruences/GaloisRepresentations/RowsAndColumnsFromLabels.lean\#L548}{IsRepOfType.isMainSummandLike}:
\[
              |\tau(y)^{\ell^u t}-1|_\ell=1
 \qquad (u\ge0,\ 1\le t\le4\,\mathrm{rank}+2).
\]
Here $\mathrm{rank}$ is either $r$ or $1$; it is at most $r$ in both cases. In detail, if the displayed difference reduced to zero, some further $\ell$-power would give $\tau(y)^{\ell^{u+v}t}=1$, contradicting Lemma~\ref{lem:scalar} and the injectivity of $\tau$. This replaces the original invocation of \texttt{not\_q\textsubscript{2}\_dvd\_mul\_pow\_mul}; it does not alter the definition of the tame order $H$ used elsewhere.

\subsection{Matrix span and scalar Frobenius}
For completeness, let $D_0$ be the inertia matrix and $F_0$ a Frobenius matrix over the local coefficient ring, with convention $F_0^{-1}D_0F_0=D_0^q$. The existing construction gives a characteristic root $\tau(z)$ and the polynomial killing relation. By Lemma~\ref{lem:erasure}, after reduction the $r$ eigenvalues $\overline{\tau(z)}^{q^i}$ are distinct. In the residue field the vectors $F_0^iv$, starting from an eigenvector $v$, are therefore independent and form a basis. The diagonal entries of $1,D_0,\ldots,D_0^{r-1}$ are a Vandermonde matrix with nonzero determinant; multiplying by the $r$ powers of $F_0$ spans every cyclic diagonal. This gives the whole matrix algebra over the residue field, and Nakayama's lemma lifts it to the coefficient ring.

Also $F_0^r$ commutes with $D_0$ because all roots are fixed by $q^r$-Frobenius; it commutes with $F_0$ automatically. In characteristic zero the same Vandermonde argument gives the full matrix algebra generated by the pair, so $F_0^r$ is scalar. The cyclotomic-character calculation of this Frobenius is unchanged. The common norm power gives $D_0^E=\tau(y)I$ from the polynomial killing relation, even though the norm is not a generator of $\F_q^\times$. These are precisely the hypotheses and conclusions of the existing \src{New/Congruences/GaloisRepresentations/RowsAndColumnsFromLabels.lean\#L510}{eventually\_span\_eq\_top\_and\_exists\_scalar\_of\_tame\_pair}; there is no need to reprove the lattice theorem in Lean.

The other inertia estimates, including determinants of $R(\gamma)-1$ at every chosen place, use exclusion of the root $1$ in the scaled type polynomial and the degree bound for its factors, not primitivity. Because the all-coordinate allowed test is unchanged, these proofs are unaffected. At each coefficient prime select the witness from the level prime different from it, exactly as before. One must not replace two level primes by one: the tame-pair argument requires the selected place to be off $\ell$.

\section{An exact count and a simpler bad-type bound}
\begin{theorem}[Witness density]\label{thm:count}
For $q\equiv1\pmod6$ and $r=3^a>1$, the number $W$ of robust polynomials of degree $r$ is
\[
                  W=\frac{q^r-1}{3r}\left(1-\frac1C\right)
                  \ \ge\ \frac{q^r}{8r}.
\]
\end{theorem}
\begin{proof}
The integer $6C$ divides $N$, because $N=3C(A-1)$ and $A$ is odd. Also $\gcd(C,6)=1$. In each block of $6C$ residue classes the Chinese remainder theorem gives precisely $2(C-1)$ exponents satisfying \eqref{eq:hexp}: two of six residues are coprime to $6$, and all but one of the residues modulo $C$ are nonzero. Thus the number of robust elements is $N(1-1/C)/3$. Every such element has degree $r$ by Lemma~\ref{lem:erasure}; Frobenius conjugacy classes have exactly $r$ elements, giving the formula. Since $C\ge5$ and $q^r-1\ge q^r/2$, the count is at least $2q^r/(15r)>q^r/(8r)$.
\end{proof}

No estimate for $\varphi(q^r-1)$ is used. In a Lean implementation even the exact CRT count can be skipped: there are $N/3$ classes coprime to $6$; at most $N/C$ of those are divisible by $C$. Since $C\ge5$, at least $2N/15$ classes remain, which alone would yield $W\ge q^r/(15r)$. The sharper $1/(8r)$ bound is preferred here, but both are much stronger than the old sufficient fraction. If a shorter proof is the overriding objective, the crude bound suffices with $15r$ replacing $8r$ below and the same admissible parameters (indeed $2(2r\log q+1)>15$ for $q\ge7,r\ge3$).

\subsection{Transfer to the regular representation}
Let $f$ be the number of chosen places above one level prime. Keep the source's uniform lower weight for each label,
\[
 w_0=\frac{1-r/q}{q^r},\qquad
 \varepsilon=\frac{1+H+2r}{q},\qquad u=\frac{1}{8r}.
\]
For $1<r<q$ and $0\le\varepsilon\le1$, the allowed labels have uniform weight at least $1-\varepsilon$, from \src{New/Spanning/MainTerm/BadTypes.lean\#L134}{one\_sub\_le\_card\_allowed\_mul}. All robust polynomials are allowed (irreducible of degree $>1$). The number of monic degree-$r$ polynomials lacking robustness is $q^r-W$. Its uniform weight is at most
\[
         (q^r-W)w_0\le 1-W/q^r\le1-u\le e^{-u}.
\]
Define the good tuples by the unchanged all-allowed condition and at least one robust coordinate. In the uniform product weight their mass is at least
\[
             (1-\varepsilon)^f-e^{-fu},
\]
since the excluded all-nonrobust tuples have weight at most $e^{-fu}$. This is a lower bound for the true Plancherel mass: every included tuple is made of labels, whose idempotents are orthogonal and have coefficient at the identity at least $w_0^f$ by the original Green degree bound.

\begin{corollary}[Bad types are rare with the new criterion]\label{cor:bad}
Let $e_q^{\rm rob}$ be the sum of tuple idempotents selected by the robust test. Then
\begin{equation}\label{eq:mass}
  \operatorname{Re}\bigl((1-e_q^{\rm rob})_1\bigr)
       \le f\frac{1+H+2r}{q}+\exp\left(-\frac{f}{8r}\right).
\end{equation}
Here the subscript $1$ denotes the coefficient at the group identity, i.e. the discarded fraction in the regular representation, not an arbitrary automorphic representation.
\end{corollary}
\begin{proof}
Subtract the preceding lower mass bound from $1$ and use Bernoulli's inequality $(1-\varepsilon)^f\ge1-f\varepsilon$.
\end{proof}

The readable source uses $u_{\rm old}=1/[2r(2r\log q+1)]$ and a sufficient condition $2r(2r\log q+1)L\le f$. We can replace it by $8rL\le f$. For the actual $q\ge7$, $r\ge3$, the old condition implies the new one; the existing parameters thus suffice without further construction. Taking $L=9\log b$ yields
\[
      f\varepsilon+e^{-f/(8r)}\le f\varepsilon+b^{-9}\le b^{-8}/2,
\]
where the last inequality is exactly the unchanged \texttt{h\_small} of \src{New/Spanning/SecondMoment/SpanningTheorem.lean\#L221}{the bad-idempotent estimate}.

The sufficient number of chosen places for the \emph{witness} part alone falls from $18r(2r\log q+1)\log b$ to $72r\log b$, an improvement by $(2r\log q+1)/4$. This is a simplification of the proof interface, not a claim that the tower degree or total number of places can be lowered by this factor. Those parameters also support the sparse places, Selmer analysis and other estimates. We keep the old values of $f$, $r$, $m$, $H$, the level primes and all cutoff scales.

\section{Preserving the rest of the argument}\label{sec:integrate}
The following dependency analysis is based on the source and guide, including direct inspection of the relevant declarations and a targeted code search for \texttt{IsPrimitivePolynomial}, \texttt{goodTuples}, \texttt{IsCommonRootPower}, and their callers. The audit is not a complete traversal of the 300,000-line development.

\begin{center}
\renewcommand{\arraystretch}{1.19}
\begin{tabular}{@{}>{\raggedright\arraybackslash}p{.28\linewidth}>{\raggedright\arraybackslash}p{.30\linewidth}>{\raggedright\arraybackslash}p{.34\linewidth}@{}}
\toprule
Source interface & What it needs & Replacement or unchanged proof \\
\midrule
Good tuple projector & Orthogonal label idempotents; allowed at every place & Same sums with robust witness instead of primitive witness. All-place hyperplane annihilation unchanged. \\
Trace comparison & At least one irreducible label giving a cuspidal test function & Robust label is irreducible. No change to the trace formula or multiplicity equality. \\
Bad type estimate & Small coefficient of $1-e_q$ & Corollary~\ref{cor:bad}; preserve $b^{-8}/2$. The existing moment transfer is reused. \\
Lift setting & A suitable witness above each level prime; scaled polynomials avoid $1$ & Replace primitive by robust; keep allowed and rank-one branches. \\
Main summand & Integral matrix span, scalar Frobenius, large-order scalar inertia & Lemmas~\ref{lem:erasure}--\ref{lem:scalar}; same theorem type \texttt{IsMainSummandLike}. \\
A, B, C and endpoint & Same uniform numerical bounds and shape of target spaces & Reinstantiate with the changed projector; unchanged inequalities, subject to build audit. \\
\bottomrule
\end{tabular}
\end{center}

The role of the trace comparison deserves emphasis. In \src{New/Spanning/Multiplicities/EqualMultiplicities.lean\#L249}{trace\_mul\_tupleIdem\_eq}, the proof destructures a primitive witness into its label property and irreducibility, discarding the primitive-root evidence. An irreducible robust label suffices for the same local cuspidal test. There is no need to identify individual forms across signatures. Likewise the code in \src{New/Congruences/GaloisRepresentations/AttachedRepresentations.lean\#L81}{isLiftSetting\_descentPair} obtains the witness and the allowed-root condition separately. The latter still proves that the scaled type polynomial never has root $1$; the former is replaced by the robust predicate. The rank-one critical-pair constructor is explicitly retained.

After these substitutions the Galois and congruence arguments receive the same \texttt{IsMainSummandLike} property. We are not modifying the Selmer bounds or replacing first-order congruences by an unproved assertion. The existing short forms, norm comparison, local losses and lower determinant heights depend on the projection and on properties of its chosen local labels. Wherever they use a general property of a good tuple (label, allowed, irreducibility, or the above lift setting), the replacement supplies that property. The remaining Lean task is to rebuild the changed projector's dependent definitions and verify that these interfaces really suffice throughout the transitive import graph.

Once that integration has been checked, the same numerical contracts give the same \texttt{star\_inequality} and hence the same one-sided analytic criterion. This is also compatible with the proposed improved spanning, signed-gain and sparse-cutoff estimates: they need the stated numerical moment or volume inputs, not the exact definition of the robust witness. No new claim is made for automorphic L-functions beyond the five original public statements.

\section{Nonduplication and implementation}
The eight proposal directories were inspected through GitHub. Their READMEs and mathematical TeX source headings were read; relevant terms (primitive polynomial/label, \texttt{IsPrimitivePolynomial}, \texttt{IsCommonRootPower}, \texttt{BadTypes}, Plancherel, \texttt{SelmerHypotheses}, \texttt{RowsAndColumnsFromLabels}, and \texttt{goodTuples}) were searched in the TeX files and available formalization plans. No matches to these interfaces were found in those texts. This establishes a targeted non-overlap assessment, not a claim of absolute novelty in the mathematical literature.
\begin{center}
\renewcommand{\arraystretch}{1.16}
\begin{tabular}{@{}p{.32\linewidth}p{.59\linewidth}@{}}
\toprule
Existing proposal & Main scope (not repeated here)\\
\midrule
\prp{centered_theta_lift_simplification}{Centered theta lifts} & Centered rank gains and adaptive basis completion.\\
\prp{cyclotomic-interpolation-and-theta-lifts}{Cyclotomic interpolation} & Selection of torsion twists and improved near-spanning.\\
\prp{dyadic_place_coloring}{Dyadic place coloring} & Sparse infinite-place construction.\\
\prp{proof_simplifications}{Proof simplifications} & Auxiliary prime, Perron inversion and spanning.\\
\prp{rh-simplifications}{RH simplifications} & Torsion interpolation, rank gain and prime selection.\\
\prp{rh_sparse_theta_proof}{Sparse theta proof} & Sparse cutoffs, local averages and twist selection.\\
\prp{shorter-theta-lift-proof}{Shorter theta-lift proof} & Character-sum/sieve arguments and spanning.\\
\prp{signed_theta_gain_simplification}{Signed theta gain} & Signed local gains and colength normalization.\\
\bottomrule
\end{tabular}
\end{center}

\paragraph{Recommended implementation sequence.}
First implement the finite cyclic-group lemma, preferably with a separate robust polynomial predicate in the local GL library. Prove its degree, count, scaling and norm lemmas. Then change the good-tuple selector and replace the count in \texttt{BadTypes.lean}. Rebuild \texttt{SpanningTheorem.lean} with the same error bound. Next adapt trace comparison and the two lift-setting constructors. Finally weaken only the scalar divisibility clause of \texttt{IsCommonRootPower}, prove the rank-one bridge by divisibility, and replace the scalar-separation invocation. Keep a regression theorem that the old primitive predicate implies the new one. The companion \texttt{LEAN\_PLAN.md} records the file-level changes and build gates.

\paragraph{Further research questions.}
\begin{enumerate}
\item Can the two-prime scalar certificate be expressed directly as the unit conclusion rather than the intermediate $6^s$ divisibility? This might make the Galois input independent of the special rank-one construction.
\item For general ranks (not powers of $3$), orbit robustness is equivalent to a two-prime condition on each quotient $(q^r-1)/(q^{r/p}-1)$ for primes $p\mid r$. Can one find a similarly dense predicate with few explicit power tests? This would require auditing every other use of the special rank in the proof.
\item Is the stronger count $W$ enough to reduce the field tower in concert with the dyadic place-coloring proposal? The improved witness bound alone is not sufficient evidence, because other errors also depend on $f$.
\item Can the exact weighted Green degree formula further relax the all-allowed test without compromising the hyperplane annihilation and the non-root condition? A large Plancherel mass alone cannot replace the exact local annihilation identity.
\end{enumerate}

\section*{Sources, verification and limits}
\addcontentsline{toc}{section}{Sources, verification and limits}
The source baseline is the ProveIt main snapshot \href{https://github.com/VladimirReshetnikov/ProveIt/tree/ee10f937fedac0955d28b7b860f4c44fab508204/Analysis/GeneralizedRiemannHypothesis}{\texttt{ee10f937fedac0955d28b7b860f4c44fab508204}}. The key source blobs and proposal directory trees are recorded in the accompanying TSV files. The principal source for the argument is the readable Lean package and the \href{https://github.com/VladimirReshetnikov/ProveIt/blob/ee10f937fedac0955d28b7b860f4c44fab508204/Analysis/GeneralizedRiemannHypothesis/docs/guide/theta_lifts_and_dirichlet_zeros.tex}{human-readable guide}, especially Section 10 (local types) and Sections 13--15 (congruences). For the original independent checks see the \href{https://math.mit.edu/~drew/shortgrh/}{Sutherland verification record} cited by the repository, and for discovery context see \href{https://alpo.ge/remark.pdf}{Alp\"oge's remarks}. This note uses the repository's reported original checks as context, not as certification of the new argument.

The exact cyclic-group regression script checks both directions of Lemma~\ref{lem:erasure} on all exponents for two small fields, the exact witness count, strict inclusion, tame invariance, and sample tests at rank nine. It uses only Python's standard library. The checks support, but do not replace, the general proofs. No Lean executable or full pinned Mathlib/project cache was available; no new declarations were kernel-checked and no repository modifications were made. A source manifest verifier is included for local reproduction. A complete proof of RH from the revised Lean sources requires successful integration and independent checking.

\end{document}
